\documentclass[11pt]{article}
\usepackage[utf8]{inputenc}
\usepackage[T1]{fontenc}
\usepackage{lmodern}
\usepackage[margin=1in]{geometry}
\usepackage{amsmath,amssymb}
\usepackage{textcomp}
\usepackage[dvipsnames]{xcolor}
\usepackage{enumitem}
\usepackage{parskip}
\usepackage{titlesec}
\usepackage[colorlinks=true,linkcolor=blue,citecolor=blue,urlcolor=blue]{hyperref}

\titleformat{\section}{\large\bfseries}{\thesection}{1em}{}
\titleformat{\subsection}{\normalsize\bfseries}{\thesubsection}{1em}{}
\titleformat{\subsubsection}{\normalsize\bfseries\itshape}{}{0em}{}

\title{\textbf{Presence Is Not Compliance:\\ Measuring When Constraint Pinning Backfires at Scale}}
\author{}
\date{}

\newcommand{\ph}{\textit{[placeholder --- fill in after experiments are run]}}

\begin{document}
\maketitle
\vspace{-2em}

\begin{quote}
\small\textit{Assembled from the six confirmed stage outputs of an \texttt{auto-research-discussion} run (seed topic T3.4), post research-taste audit (\textcolor{ForestGreen}{\textbf{PASS}}). Status: \textbf{specification / paper draft --- no experiments have been run.} All numbers below are pre-registered design targets, not results.}
\end{quote}

\begin{abstract}
LLM agents are increasingly asked to honor large standing constraint sets --- safety rules, user preferences, organizational policy --- and the default engineering fix for making sure such an agent ``sees'' every rule is pinning: repeating the full constraint set verbatim in every turn. We show this default rests on an unexamined asymmetry: the token cost of pinning a constraint scales linearly with how many constraints there are, while the attention budget that pays for actually complying with each one does not --- so the practice that guarantees a rule is present in context is not guaranteed to guarantee compliance with it, and nothing in current practice locates the point, if any, where this gap becomes costly. We build a constraint-count sweep harness that deconfounds constraint count from raw prompt length and from applicability density, instrument full pinning's token overhead directly against task success to locate (or rule out) a breakpoint, evaluate a family of budget-capped, step-conditioned selective-pinning policies --- including an applicability-aware selector tested against similarity-based ones --- against a random-budget floor, and gate every violation-rate curve behind a human-calibrated, dual-mode violation detector. We report where, if anywhere, full pinning's token overhead begins to erode task success as the constraint count scales from single digits to several hundred, and whether a bounded-budget selection policy can hold the line on violations at a fraction of that cost.
\end{abstract}

\section{Introduction}

Production LLM agents are now asked to honor hundreds of standing rules at once. Airline and retail policy agents, enterprise copilots, and personal assistants are each expected to respect a growing accumulation of safety constraints, organizational policy, and user preferences --- not one or two, but potentially hundreds of standing rules that must all hold simultaneously across a multi-turn task. The default engineering answer to ``how do we make sure the agent doesn't forget a rule'' is pinning: paste every constraint, verbatim, into the context window of every single turn. It is cheap to implement, easy to reason about, and treated in practice as a safety default rather than an open empirical question.

Yet nothing establishes whether the default fix for this --- pinning --- keeps working as the rule count grows. Pinning purchases a constraint's \emph{presence} in context at a token price that scales linearly with how many constraints there are, but compliance with any one constraint is paid for out of a fixed per-turn attention budget that every other co-resident constraint dilutes. These two currencies do not have to move together: a rule can be perfectly present in every turn's context and still be violated, if the sheer number of co-resident rules degrades the model's fidelity to any single one of them, or if the growing constraint block itself starts crowding out the tokens the task needs to succeed. Whether either failure mode actually occurs for standing normative constraints in a multi-turn agent setting --- and if so, at what constraint count, at what token price, and with what consequence for violations --- has never been measured jointly.

The nearest existing work approaches this gap from adjacent but distinct angles. Work on multi-constraint instruction following, such as RECAST, trains models and mines verification rewards to improve how many simultaneous requirements a model \emph{satisfies}, closest in domain to this question but measuring satisfaction, not the token cost of carrying the constraint set that must be satisfied. Work on controlled cognitive-load benchmarking, such as CogniLoad, is closest in method --- it already decouples distractor density, sequence length, and task complexity as independently tunable factors so that degradation can be attributed to one dimension rather than conflated across several --- but has only ever been applied to factual and reasoning material, never to a standing set of normative rules whose violation, not incorrectness, is the outcome of interest. A concurrent late-2025--2026 cluster --- learnable context management for long-horizon agents, bounded-memory testbeds, hindsight-replay training for multi-constraint following --- budgets what enters an agent's window or raises how many constraints get satisfied, but likewise prices neither the token cost of carrying a standing normative set nor its violation consequence at scale (Related Work, concurrent-work paragraph). Neither the nearest prior work nor this concurrent cluster establishes whether, or where, a growing constraint count itself --- as opposed to a longer prompt, or a harder task --- degrades an agent's compliance.

Turning this gap into a measurable question surfaces its own difficulties. Any observed decline in task success as the constraint count grows is naturally entangled with the fact that pinning more constraints also means a longer prompt, so an effect attributed to ``constraint count'' could just as easily be an artifact of raw prompt length unless the two are experimentally separated. Which constraints actually bind at a given step is a state- and task-dependent question that must be re-decided at every turn, not a static property a similarity search can be assumed to capture, so a selection policy built on semantic relevance alone might be selecting the wrong subset for reasons that have nothing to do with meaning. The three quantities that matter --- violation rate, token overhead, and task success --- do not collapse into one honest scalar, because a safety rule and a formatting preference do not price a violation the same way. And because exhaustive human checking of hundreds of constraints across many episodes is infeasible, every one of these measurements is only as trustworthy as whatever automatic detector produced it, which means detector error is not a side issue but a confound on every downstream claim.

Our approach turns each of these obstacles into a specific design choice rather than a hedge. A factorial sweep harness holds constraint count, raw prompt length, and applicability density independently controllable --- including a length-matched inert-padding arm that replaces constraints with equal-length non-normative text --- so that a decline attributable to constraint count can be told apart from one attributable to length alone. Full pinning's token overhead is instrumented directly against task success across the sweep, and a segmented regression with a formal breakpoint test locates a knee in the trade-off if one exists, rather than assuming one. A family of five budget-capped selectors, re-decided every turn against ground-truth applicability labels, tests whether an applicability-aware selector actually outperforms selectors built only on semantic similarity to the current step. And every violation-rate curve is gated behind a dual-mode detector --- exact programmatic checkers for rule-based constraints, an LLM judge for semantic ones --- calibrated against human annotation before any curve is read as fact.

We evaluate this design across a constraint-count sweep spanning single digits to several hundred simultaneously active constraints, drawn from real domain-policy and verifiable-instruction corpora, comparing full pinning against no pinning, a family of selective-pinning policies, and a random-budget floor, on tool-use agent tasks with programmatically checkable end states.

In summary, we contribute:
\begin{itemize}[leftmargin=1.4em]
\item An identification and formalization of the asymmetry underlying constraint pinning as a governance default: the token cost of guaranteeing a rule is \emph{seen} scales linearly with rule count, while the attention budget that pays for compliance to each rule does not, so presence and compliance can silently decouple as the standing rule count grows.
\item A factorial constraint-count sweep harness that deconfounds constraint count from raw prompt length and applicability density, letting a compliance decline be attributed to the right cause rather than assumed.
\item A budget-capped, step-conditioned selective-pinning policy family --- including a direct test of whether normative applicability reduces to semantic similarity --- evaluated against a capability-matched random-budget floor, not benchmarked as a system meant to beat a leaderboard.
\item An empirical account of whether, and at what constraint count, full pinning's token overhead begins to measurably erode task success, gated throughout by a human-calibrated violation detector, and what this implies for pinning as a default governance mechanism in production LLM agents.
\end{itemize}

\section{Problem Definition}

\noindent\textbf{Formal setup.} Consider an episode $e = (\tau, C, \pi, m, s)$: a task $\tau$ drawn from a task suite, a constraint set $C$ with $|C| = n$ drawn from a constraint pool, a pinning policy $\pi$, a model $m$, and a random seed $s$. The episode unfolds over turns $t = 1,\dots,T_e$. Each constraint $c \in C$ carries a ground-truth applicability label $\mathrm{app}(c,e,t) \in \{0,1\}$ --- whether $c$ actually binds the agent's action at turn $t$ --- and a verification procedure that yields a detector verdict $\mathrm{viol}(c,e) \in \{0,1\}$ on turns where $c$ is applicable. The pinning policy $\pi$ governs which subset of $C$ is resident in the agent's context at each turn: $\pi_{\text{full}}$ repeats all $n$ constraints verbatim every turn; $\pi_{\text{none}}$ states the constraint set once, at $t=1$; $\pi_k^{\sigma}$ selects, under selector $\sigma$ and a fixed token budget $k$, a bounded subset to keep resident at each turn. Three quantities are tracked as functions of $(n, \pi, m)$: a violation rate $V$, the fraction of applicable (constraint, episode) pairs that are violated; an overhead ratio $O$, the mean per-turn fraction of the prompt occupied by the constraint block; and a task success rate $S$, the fraction of episodes passing a programmatic end-state check. The object of study is whether $\mathrm{logit}(S)$ as a function of $\log n$ under $\pi_{\text{full}}$ has a detectable breakpoint --- and, if so, whether a budget-capped selective policy can hold $V$ near $\pi_{\text{full}}$'s level at a fraction of $\pi_{\text{full}}$'s $O$.

\noindent\textbf{Why this is hard.} Nothing in the pinning practice itself imposes a ceiling on $n$ --- the token cost of $\pi_{\text{full}}$ scales linearly with no built-in circuit breaker, so whether and where $S$ actually bends as $n$ grows is an empirical fact to locate, not a design parameter to read off; a segmented model with a formal breakpoint test is needed to turn this from an eyeballed trend into a falsifiable claim. Any observed decline as $n$ grows is also confounded with the fact that a larger $C$ means a longer prompt, so isolating a genuine constraint-count effect from a raw-length effect requires holding length fixed while varying only constraint content, and holding applicability density fixed while varying count. Which constraints actually bind at a given turn --- $\mathrm{app}(c,e,t)$ --- is a state- and task-dependent normative-matching question, not a semantic-similarity question, and no existing selection method is validated against ground-truth applicability rather than textual relevance. There is no single scalar that prices $V$, $O$, and $S$ against one another correctly, since different constraint classes (a safety rule versus a formatting preference) do not price a violation equally. And because exhaustive human checking of hundreds of constraints across many multi-turn episodes is infeasible, every measured $V$ is only as trustworthy as the detector that produced $\mathrm{viol}(c,e)$, so detector error must be bounded empirically before any curve derived from it is interpreted as fact.

\section{Method Overview}

At a conceptual level, the study is a single measurement pipeline with five cooperating parts, four of which produce a falsifiable claim and one of which is deliberately descriptive. Constraint pools and agent task episodes enter a factorial sweep harness that independently controls constraint count, raw prompt length, and applicability density; full pinning's token cost is instrumented directly against task success across that sweep to locate a breakpoint if one exists; a family of budget-capped selectors, evaluated within the same harness, turns the sweep into a test of whether a bounded selection policy can hold the line on violations; a dual-mode detector gates every violation reading behind a human-calibrated agreement threshold; and a trade-off tracker renders the resulting three-axis surface without collapsing it into a single number.

The \textbf{deconfounded sweep harness} takes the constraint pool and task suite and produces episodes in which constraint count, prompt length, and applicability density are each independently manipulable --- including a length-matched inert-padding arm --- so that a decline in compliance can be attributed to the right cause. The \textbf{overhead instrumentation and breakpoint test} takes the resulting per-turn token counts and task-success outcomes under full pinning and fits a segmented model that either locates a breakpoint in $n$ or reports, as a valid finding, that none exists in the tested range. The \textbf{selective-pinning policy family} takes the same harness under a fixed token budget and evaluates five selectors --- including one built on ground-truth applicability rather than textual similarity --- to test whether a bounded policy can approximate full pinning's violation rate at a fraction of its cost. The \textbf{dual-mode violation detector} takes every episode's transcript and produces a calibrated violation reading, gating programmatic and semantic constraints differently and validating both against human annotation before any curve is trusted. The \textbf{trade-off tracker} takes the full $(n, \pi, m)$ grid of violation rate, overhead, and success and renders it as a Pareto surface with per-constraint-class breakdowns, reporting the characterized frontier rather than resolving it to one preferred operating point.

\section{Methodology}

\subsection{Deconfounded Sweep Harness}
\textbf{Challenge.} Any observed decline in task success as constraint count grows is naturally confounded with the fact that a larger constraint set also means a longer, more attention-dividing prompt --- a claim about ``count'' could always really be a claim about ``length'' or ``how much of the set actually binds'' instead.

\textbf{Insight.} A count effect, a length effect, and an applicability-density effect make the same prediction (success declines as more text is pinned) unless something holds two of the three fixed while varying only the third --- deconfounding therefore requires a factorial design, not a bigger sample of the naturally-correlated raw condition.

\textbf{Method.} Episodes are instantiated with constraint count $n$, applicability density, and prompt length as independently controllable factors. A length-matched inert-padding arm replaces constraints with equal-length non-normative filler text, holding token count fixed while removing all normative content, so that a decline reproduced by padding alone is attributable to length, and a decline that survives only under real constraint content is attributable to constraint count itself. Ground-truth applicability labels $\mathrm{app}(c,e,t)$ are authored at pool-construction time so that density can be independently varied at fixed $n$. Because these labels are load-bearing three times over --- they set the density factor, define the violation-rate denominator, and feed the oracle selector --- a stratified sample of them is independently re-annotated by a second annotator, with agreement reported alongside RQ4's detector calibration, so that label noise cannot silently distort the quantities built on them.

\subsection{Overhead Instrumentation and Breakpoint Test}
\textbf{Challenge.} Nothing in the pinning practice imposes a ceiling on constraint count, so whether --- and where --- the linear growth in token cost actually starts eroding task success is an empirical fact that has to be located, not a design parameter that can be read off in advance.

\textbf{Insight.} A breakpoint claim needs a formal test with a stated significance threshold and a confidence interval, not an eyeballed inflection in a plotted trend, or ``it backfires past some point'' is unfalsifiable by construction.

\textbf{Method.} Under full pinning, per-turn constraint-block token counts (absolute and as a fraction of the prompt) are logged alongside task-success outcomes across the constraint-count sweep. A segmented (piecewise-linear) regression of $\mathrm{logit}(S)$ on $\log n$ is fit, and a Davies test checks for a breakpoint; the breakpoint's location, if detected, is reported with a bootstrap confidence interval. Violation rate under full pinning is tracked in parallel across the same sweep, so that a rising violation rate \emph{without} a corresponding drop in the constraints' presence can be read as compliance decoupling from presence, rather than assumed.

\subsection{Selective-Pinning Policy Family}
\textbf{Challenge.} Which constraints actually bind at a given turn is a state- and task-dependent normative question --- ``does rule $r$ apply to the action about to be taken'' --- not a question that textual or semantic similarity to the current step is guaranteed to answer correctly, and no prior selection method has been validated against ground-truth applicability rather than relevance.

\textbf{Insight.} The only way to test whether normative applicability reduces to semantic similarity is to build a selector that is deliberately not similarity-based and compare it, at the same token budget, against selectors that are --- a difference in outcome between the two directly answers the question that similarity-only selection methods leave untested.

\textbf{Method.} Under a fixed token budget $k$, five selectors are implemented and re-run every turn: a random-$k$ floor, a static top-$k$ policy ranked by pilot-measured violation frequency, an embedding-similarity selector, a marginal-information-gain selector (adapting the scoring rule used for evidence compression), and an applicability classifier scored directly against the ground-truth labels from the sweep harness. Because this last selector consumes applicability labels that would not be available at deployment time, it is read as an \emph{oracle upper bound} on applicability-aware selection: the applicability-versus-similarity contrast is an existence test of how much headroom lies above semantic relevance, not a claim that a deployable classifier already attains that headroom. Each selector's resulting violation rate and overhead are measured at matched budgets, so that any gap between the applicability classifier and the similarity-based selectors is attributable to the selection criterion itself, not to budget or task differences.

\subsection{Dual-Mode Violation Detector}
\textbf{Challenge.} Exhaustive human checking of hundreds of constraints across many multi-turn episodes is infeasible, yet every violation-rate curve in this study is only as trustworthy as whatever automatic procedure produced it --- an uncalibrated detector could manufacture or mask the very trade-off being characterized.

\textbf{Insight.} Detector error is a confound on every downstream claim, not a separate side-question to address after the fact --- gating on a pre-registered human-agreement threshold before any curve is interpreted is what prevents a noisy detector from silently deciding the paper's findings.

\textbf{Method.} Rule-based constraints are checked with exact programmatic verifiers; semantic constraints are checked with an LLM judge. Both are calibrated against a human-annotated, stratified sample (by constraint count, policy, and constraint class) before any curve derived from the detector is read as a finding. If the pre-registered agreement threshold is not met, headline claims are restated on the programmatically-checkable subset of the pool alone, with semantic-constraint results reported separately and flagged as detector-error-bounded rather than withdrawn silently.

\subsection{Trade-off Tracker}
\textbf{Challenge.} Violation rate, token overhead, and task success do not collapse into one honest scalar objective, because different constraint classes do not price a violation the same way, and no data-driven procedure can justify a single weighting across them.

\textbf{Insight.} Reporting the full three-axis surface, broken down by constraint class, is the honest characterization; forcing a single number would bake in an arbitrary value judgment the design itself cannot defend.

\textbf{Method.} The full grid of (constraint count, policy, model, seed) outcomes is logged and rendered as a Pareto surface with per-constraint-class breakdowns. This component deliberately produces no falsifiable hypothesis of its own: it describes the resulting trade-off surface rather than resolving it to a single recommended operating point, and that description --- not a claim about where the ``right'' point on the frontier sits --- is its contribution to this study.

\section{Experiments}

\subsection{Experimental Setup}

\textbf{Datasets.} $\tau$-bench (airline/retail policy-following tool-use tasks) supplies the primary task suite, with task success verified via database-state checks. IFEval and FollowBench supply rule-based, programmatically verifiable constraint templates, mined and scaled (in the style of RECAST's constraint-mining pipeline) to build the constraint pool up to $n=400$. PrefEval supplies the semantic-constraint templates that make up the LLM-judged remainder of the pool, which the detector-validity research question specifically needs to stress-test.

\textbf{Baselines} (floor and sanity-check comparators, not competitors to beat --- this is a measurement study, and no research question here is framed as ``our procedure beats baseline X''). Full pinning is the stress-tested default every condition is compared against. No pinning (constraints stated once, at episode start) is the necessary control establishing that the underlying problem exists at all: if violation rate under no pinning does not rise with constraint count, constraint retention is a non-problem and pinning of any kind needs no further study. A random-$k$ selector at the same token budget is the capability-matched floor that any claim about applicability-aware selection must clear. A length-matched inert-padding arm is the rival-explanation control isolating constraint count from raw prompt length.

\textbf{Metrics.} Violation rate and its mixed-effects trend slope; token overhead (absolute and as a fraction of the prompt); task success rate and its segmented-regression breakpoint (if any); and detector agreement (Cohen's $\kappa$ against human annotation).

\textbf{Measurement cost.} The scale-up grid totals 7 constraint-count levels $\times$ 6 policies $\times$ 3 models $\times$ 450 episodes per cell $=$ 56{,}700 multi-turn tool-use episodes, on top of the pilot --- and the full-pinning cells dominate the token bill, since their constraint blocks grow linearly in $n$ up to $n=400$ in every turn of every episode. A study whose thesis is that linear-in-$n$ token cost is a problem owes its own cost accounting: per-condition token spend is logged as a first-class output (the same instrumentation that produces the overhead ratio $O$ prices the experiment itself and is released with the artifact), the invariant pinned constraint block is served through prompt-prefix caching wherever the serving stack supports it --- with cached and uncached token counts reported separately, since caching changes the price of presence but not the attention-budget load under study --- and the scale-up is gated on the pilot, so the largest cells are run once, not iterated.

\textbf{Research questions:}
\begin{enumerate}[leftmargin=1.6em]
\item RQ1 --- Does full pinning's task success have a detectable breakpoint as constraint count grows, and does violation rate rise even while every constraint remains resident?
\item RQ2 --- Can a budget-capped selective-pinning policy hold violation rate near full pinning's level at a fraction of its overhead, and does an applicability-aware selector beat similarity-based ones?
\item RQ3 --- Is the observed decline attributable to constraint count itself, as opposed to raw prompt length or applicability density alone?
\item RQ4 --- Is the dual-mode violation detector reliable enough that the measured curves reflect agent behavior rather than detector error?
\end{enumerate}

\subsubsection{RQ1: Does full pinning's task success have a detectable breakpoint, and does violation rate rise even while every constraint remains resident?}
\textbf{Setup.} Pilot stage: 2 backbones $\times$ 2 conditions (full pinning / no pinning) $\times$ 50 pairs $\times$ 20 episodes on the rule-based subset. Scale-up: 7 constraint-count levels $\{4, 10, 25, 50, 100, 200, 400\}$ $\times$ 6 policies $\times$ 3 models $\times$ 150 tasks $\times$ 3 seeds (450 episodes/cell). Mixed-effects logistic regression of episode success on $\log n$, task as random effect; segmented regression with a Davies test for a breakpoint, $\alpha = 0.05$, Holm-corrected across models; violation rate compared via cluster-robust standard errors at the episode level.

\textbf{Sample-size rationale.} 450 episodes/cell is sized to detect a 10-percentage-point success difference at 80\% power, $\alpha = 0.05$. Any cell falling below 388 episodes/arm after dropped-row exclusion is pre-declared INCONCLUSIVE for that cell --- reported with its achieved sample size, never scored as a null result and never silently dropped; a primary hypothesis whose required cells fall short is itself declared INCONCLUSIVE, with the claim narrowed to the powered subset. This budget prices the task-success contrast only: the 3-percentage-point violation-rate contrast here, and RQ2's 0.02 non-inferiority margin, are smaller effects than the success difference the cell size was computed for, and are powered at the (constraint, episode)-pair level --- each episode contributes one observation per applicable constraint, with cluster-robust standard errors absorbing intra-episode correlation --- not at the episode level. If the pilot's measured base rates and intra-episode correlation imply less than 80\% power for a violation-rate margin at the pair level, the affected margin is relaxed to a pre-registered fallback of 0.03 absolute and reported as relaxed, under the same never-silently-dropped rule.

\textbf{Expected outcome.} Task success under full pinning has a detectable breakpoint (Davies test $p < 0.05$, 95\% bootstrap CI inside $[25, 400]$), with success at $n=400$ at least 10 percentage points below success at $n=10$ in at least 2 of 3 tested models; violation rate at $n=400$ exceeds violation rate at $n=10$ by at least 3 percentage points (absolute) in at least 2 of 3 models, even though every constraint remains in context at every turn.

\textbf{Falsification condition.} Refuted (no backfire) if, across the full tested range up to $n=400$, full-pinning task success shows no statistically significant decline --- the segmented regression finds no breakpoint and the coefficient on $\log n$ is non-negative --- in which case the honest conclusion is that exhaustive pinning's cost, within this range, is economic rather than behavioral. Separately refuted on the compliance-decoupling claim if violation rate at $n=400$ is statistically indistinguishable from violation rate at $n=10$ despite the token-overhead ratio having grown substantially --- that would mean presence and compliance track each other even at scale.

\textbf{Hypothesis (plain language).} Guaranteeing a rule is present in context, via full pinning, does not guarantee compliance with it, and this gap grows large enough to measurably reduce task success once the standing rule count reaches the low-to-mid hundreds.

\textbf{Metrics.} Segmented-regression breakpoint location and 95\% CI; per-cell success at $n=10$ versus $n=400$; per-cell violation rate at $n=10$ versus $n=400$; overhead ratio logged in parallel.

\textbf{Result.} \ph

\subsubsection{RQ2: Can a budget-capped selective-pinning policy hold violation rate near full pinning's level at a fraction of its overhead, and does an applicability-aware selector beat similarity-based ones?}
\textbf{Setup.} At constraint counts $n \geq 100$, all five selectors (random-$k$, static top-$k$, embedding similarity, marginal-information-gain, applicability classifier --- the last consuming ground-truth labels and read as an oracle upper bound, per the Methodology) are run at a token budget fixed to 25\% of full pinning's constraint tokens at the matching $n$. Non-inferiority of each selector's violation rate against full pinning is tested via a one-sided 95\% confidence interval at a margin of 0.02 absolute; the applicability classifier is compared against the best similarity-based selector via a Holm-corrected pairwise test.

\textbf{Expected outcome.} At $n \geq 100$, the best-performing selector achieves a violation rate non-inferior to full pinning within 0.02 absolute while using at most 30\% of full pinning's overhead; the applicability-classifier selector achieves a violation rate at least 0.03 (absolute) lower than the best similarity-based selector at the same budget.

\textbf{Falsification condition.} Refuted on viability if no selector, at any tested budget, clears the non-inferiority margin against full pinning --- the honest conclusion in that case is that for normative content, unlike evidence retrieval, no bounded selection is currently safe, and full pinning remains the only reliable policy below the breakpoint identified in RQ1. Refuted on the applicability-versus-similarity claim if the applicability classifier and the best similarity-based selector are statistically equivalent (within $\pm0.02$, under a two-one-sided-tests (TOST) equivalence test) --- that would mean normative applicability is, for this constraint pool, empirically reducible to semantic relevance, and cheaper similarity-based selection suffices.

\textbf{Hypothesis (plain language).} A bounded, step-conditioned selection policy can approximate full pinning's compliance at a fraction of its token cost, and doing so requires reasoning about which rules actually apply rather than which rules merely sound relevant to the current step --- headroom the oracle contrast can establish exists, but that a deployable selector would still have to realize without ground-truth labels.

\textbf{Metrics.} Per-selector violation rate and overhead ratio at matched budgets; non-inferiority margin test result against full pinning; pairwise applicability-classifier-versus-similarity contrast.

\textbf{Result.} \ph

\subsubsection{RQ3: Is the observed decline attributable to constraint count itself, as opposed to raw prompt length or applicability density alone?}
\textbf{Setup.} At $n=400$, a length-matched inert-padding arm (identical token count and position, no normative content) is compared against the real-constraint full-pinning arm on task-success decline relative to $n=4$. At fixed $n=100$, applicability density is independently varied at three levels, holding $n$ constant, and violation rate is compared across density levels.

\textbf{Expected outcome.} Length-matched inert padding reproduces at most 50\% of the full-pinning task-success decline observed at $n=400$ (statistically distinguishable at $\alpha=0.05$); at fixed $n=100$, high applicability density raises violation rate by at least 0.03 (absolute, $p<0.05$) relative to low density.

\textbf{Falsification condition.} Refuted (pure length artifact) if inert padding reproduces the full decline observed under real constraints --- that would mean the ``constraint count'' effect is actually a raw-prompt-length effect, and constraint-count guidance collapses into existing long-context guidance with nothing paper-specific to add. Separately, if applicability density shows no measurable effect on violation rate at fixed $n$, density is dropped from the harness as a non-load-bearing factor and this is reported explicitly rather than silently dropped.

\textbf{Hypothesis (plain language).} Constraint-count load is a distinct degradation dimension from raw prompt length --- pinning more rules is not simply a special case of ``the prompt got longer.''

\textbf{Metrics.} Task-success decline magnitude, real-constraint arm versus padding arm, with a difference-of-declines significance test; violation rate by applicability-density level at fixed $n$.

\textbf{Result.} \ph

\subsubsection{RQ4: Is the dual-mode violation detector reliable enough that the measured curves reflect agent behavior rather than detector error?}
\textbf{Setup.} A 300-episode sample, stratified by constraint count, policy, and constraint class, is independently annotated by two human raters. Cohen's $\kappa$ is computed between the detector's verdicts and human annotation, overall and per constraint class, with a bootstrap confidence interval, and the per-class confusion structure is inspected to confirm disagreements concentrate between adjacent constraint sub-types rather than collapsing one class into another.

\textbf{Expected outcome.} The dual-mode detector reaches Cohen's $\kappa \geq 0.80$ against human annotation overall, with programmatic checkers reaching at least 0.95 raw agreement, after at most one calibration iteration.

\textbf{Falsification condition.} Refuted if $\kappa$ remains below 0.80 after one calibration iteration --- in which case semantic-constraint curves are withdrawn from the headline claims (RQ1--RQ3), which are re-stated on the programmatically-checkable subset of the constraint pool alone, with LLM-judged results reported separately and flagged with judge-error bars rather than presented at equal confidence. The same withdrawal applies per constraint class if the confusion-structure check finds a systematic (non-adjacent) collapse of one class into another, even when aggregate $\kappa$ clears 0.80.

\textbf{Hypothesis (plain language).} The measured violation-rate curves reflect what the agents actually did, not an artifact of an unreliable detector --- and this needs to be checked before, not after, the headline curves are read.

\textbf{Metrics.} Overall and per-constraint-class Cohen's $\kappa$ with bootstrap CI; programmatic-checker raw agreement rate; per-class confusion-rate inspection.

\textbf{Result.} \ph

\vspace{0.5em}
\noindent\textbf{Further evaluation.} With more time, a natural extension adds a distractor-conversation dimension --- constraints buried inside long, irrelevant dialogue history rather than deliberately pinned and always resident, closer to the retention setting PrefEval studies for a single preference. This is explicitly out of scope here: every constraint in this design is, by construction, pinned or selectively pinned, never merely present somewhere in an unstructured history, so a distractor-conversation study would need its own research question and its own dataset rather than an extension of the sweep above.

\section{Related Work}

The closest prior work to this paper is RECAST and CogniLoad, and it is worth stating the delta plainly before surveying the wider landscape: RECAST is closest in domain --- it already documents that models degrade substantially as the number of simultaneously imposed requirements grows --- and CogniLoad is closest in method --- it already decouples distractor density, sequence length, and task complexity as independently tunable factors so that degradation attributes to one dimension rather than several conflated together. But RECAST trains models to better \emph{satisfy} a fixed set of requirements, never measuring the representation cost of carrying that set across a multi-turn trajectory, and CogniLoad's factorial design has only ever been applied to factual and reasoning material, never to a standing set of normative rules whose violation, rather than incorrectness, is the outcome that matters. Closing exactly this gap --- a controlled, deconfounded sweep over the count of pinned normative constraints, jointly against violation rate, token overhead, and task success --- is precisely what this paper is built to do.

\textbf{Multi-constraint instruction following.} RECAST mines verifiable constraints from instruction-response corpora, splits them into rule-based and model-based categories, and trains with constraint-level verification rewards to improve how many simultaneous requirements a model satisfies, motivated by evidence that models degrade substantially as requirement count grows. This line of work treats the constraint set as free to represent --- it optimizes satisfaction of whatever constraints are given, never the token cost of keeping a growing constraint set resident in an agent's context across a multi-turn trajectory, which is the quantity this paper measures instead.

\textbf{Standing-item retention and forgetting under long context.} PrefEval tests whether a single user preference, stated once, survives being buried inside a long, unrelated conversation; ICF-Bench tests whether a model can suppress one revoked or superseded instruction once new context supersedes it. Both isolate a single standing item against a distractor-filled history and report compliance only. Neither scales to simultaneously binding many constraints at once, and neither reports the token-overhead cost of whatever retention mechanism --- including simply re-stating the item --- might be used to keep it salient, which this paper's overhead instrumentation is specifically built to measure.

\textbf{Agent context and history management.} Sculptor equips an agent with tools to fragment, hide, and restore spans of its own task history, curating a low-interference workspace over what has already happened; AgentFold learns a folding policy that maintains multi-scale summaries of older sub-trajectories alongside a full-fidelity latest step. Both actively manage the \emph{task history} a session accumulates, but both treat the standing instruction or policy block as fixed background, never reporting constraint-violation rate as an outcome --- they optimize what enters the workspace from past interactions, not which of hundreds of standing rules deserve context residence at all, which is this paper's actual object of study.

\textbf{Concurrent bounded-context and constraint-scaling work (late 2025--2026).} A fast-moving concurrent cluster now approaches individual axes of this design's estimand, and the delta is worth conceding precisely. Context as a Tool (CAT; arXiv:2512.22087) makes context management itself a learnable tool for long-horizon software-engineering agents, deciding under a token budget what enters the window; AgenticSTS (arXiv:2607.02255) instruments a bounded-memory contract that decides which typed experience slots occupy a finite per-decision window and ablates those slots against task win-rate; and hindsight-replay RL for multi-constraint instruction following (HiR; arXiv:2512.23457) trains models against an explicit constraint set $C = \{c_1, \dots, c_n\}$ to raise how many constraints are satisfied simultaneously. Each of these occupies exactly one axis of the present design --- budgeted context residency for \emph{task history or experience} (CAT, AgenticSTS), or \emph{satisfaction} of a multi-constraint set at training time (HiR) --- and none measures what this study is built to measure: violation rate, token overhead, and task success jointly, as the count of standing \emph{normative} constraints sweeps from single digits to several hundred, with count deconfounded from raw prompt length and applicability density. The residency decisions studied in that cluster concern what an agent has done or experienced; the residency decision studied here concerns what an agent is obliged to obey, whose violation --- not omission --- is the outcome of record.

\textbf{Context compression and controlled load benchmarking.} COMI compresses context under a fixed token budget by scoring units on marginal information gain --- relevance to a query minus redundancy with what is already kept --- validated for evidence and summarization retrieval; the same marginal-gain scoring rule is reused here as one arm of the selective-pinning policy family, but its validation never extends to normative content, where ``relevance'' means whether a rule applies to the current action rather than semantic similarity to a query, and where compressing a rule's wording risks destroying its enforceability. CogniLoad, discussed above as the nearest methodological neighbor, and Lost in the Middle, which shows that a single retrieval target can go unused when buried in the middle of a long context, both establish that presence in context does not guarantee use --- but neither quantifies how compliance degrades as the number of simultaneously binding, normative items grows, nor the overhead cost of keeping them all resident, which this paper's sweep is designed to close.

\section{Conclusion}

This work reframes a governance default --- that pinning every standing constraint verbatim is a safe, if wasteful, choice --- as an empirical question with a precise identifiability structure, rather than treating verbatim repetition as self-evidently sufficient. The central methodological wager is that presence and compliance are only distinguishable once constraint count is deconfounded from raw prompt length and applicability density, and once a breakpoint claim is backed by a formal segmented-regression test rather than an eyeballed trend; weaker designs risk mistaking a length artifact for a constraint-count effect, or an unmeasured detector's noise for a real behavioral trend. Whichever outcome the sweep produces, the finding is decision-relevant: a detected breakpoint would mean pinning needs a companion selection mechanism well before the constraint count reaches production scale; no detected breakpoint within the tested range would mean the cost of exhaustive pinning is currently economic rather than behavioral, and practitioners can trade token cost for simplicity without a compliance penalty. Two limitations bound this claim from the outset and are natural next steps rather than afterthoughts: the trade-off surface across violation rate, overhead, and success is characterized, not resolved to a single recommended operating point, since no data-driven procedure can justify one weighting across constraint classes on its own; and the study is scoped to a fixed, bounded standing constraint set within a single episode, leaving constraints that are added, revised, or revoked mid-task --- and constraint sets that persist across sessions rather than within one --- for follow-up work.

\vspace{1em}
\hrule
\vspace{1em}

\section*{Pre-handoff checklist}
\begin{itemize}[leftmargin=1.4em]
\item \textbf{Organization/navigation} --- narrative flows problem $\to$ method $\to$ experiments $\to$ related work without forward references to undefined terms.
\item \textbf{Terminology discipline} --- ``pinning,'' ``constraint pool,'' ``applicability,'' ``overhead ratio,'' ``trade-off surface'' each defined at first use in Problem Definition, used consistently thereafter; no internal stage IDs (\texttt{sp\_*}, \texttt{m\_*}, \texttt{rq\_*}) leaked into prose --- checked directly against the discussion pipeline log.
\item \textbf{Notation with worked example} --- the formal setup section defines $(e, \tau, C, \pi, m, s, n, t, \mathrm{app}, \mathrm{viol}, V, O, S)$ and the Problem Definition/Methodology sections consistently reuse this notation; no separate worked numeric example is included since no numbers exist yet (would require fabricating results).
\item \textbf{RQ framing matches outcomes} --- all 4 RQs are graded/quantitative (breakpoint locations, non-inferiority margins, agreement rates), each answered in its own subsection with an explicit falsifiable expected outcome \emph{and} an explicit falsification condition, not a binary yes/no.
\item \textbf{No typesetting artifacts} --- only explicit placeholder markers in Results, per the drafting convention this is a pre-registered spec, not a report.
\item \textbf{Artifact plan} --- the constraint-count sweep harness, the labeled 400-constraint pool with per-turn applicability ground truth, and the calibrated dual-mode detector are the natural release artifacts; not yet built, so not claimed as available.
\end{itemize}

\section*{Hand-off}
This is a draft, not a finished paper --- experiments have not run. Natural next steps: prose polish and venue-specific LaTeX formatting once a target venue is chosen, or revising an earlier stage if piloting surfaces a problem (e.g., if the detector's calibration in RQ4 is much weaker than expected).

\end{document}
